\section{政策启示}

\begin{frame}{政策建议：从“做大总量”走向“精准施策”}
  \begin{enumerate}
    \item \hlb{持续完善绿色金融体系}：扩大绿色信贷、绿色债券、绿色保险、绿色基金等供给规模。
    \item \hlb{实施差异化区域政策}：西部继续聚焦清洁能源与节能技改，东中部强化绿色技术创新与产业升级。
    \item \hlb{强化政策协同}：推动绿色金融与产业政策、环境规制、碳交易体系联动。
    \item \hlb{完善基础设施}：统一绿色标准、建设信息共享平台、加强绿色金融人才培养。
  \end{enumerate}
\end{frame}

\begin{frame}{实践价值：对政府、金融机构与企业分别意味着什么}
  \begin{columns}[T,onlytextwidth]
    \column{0.32\textwidth}
    \begin{block}{政府}
      更需要看 \hl{政策效果}，而非只看工具投放规模。
    \end{block}

    \column{0.32\textwidth}
    \begin{exampleblock}{金融机构}
      需要降低“漂绿”风险，把资源投向 \hlb{真正有减排效果} 的项目。
    \end{exampleblock}

    \column{0.32\textwidth}
    \begin{alertblock}{企业}
      绿色融资不只是拿资金，更是推动低碳转型和竞争力重塑的契机。
    \end{alertblock}
  \end{columns}
\end{frame}
